% Template for PLoS
% Version 3.0 December 2014
%
% To compile to pdf, run:
% latex plos.template
% bibtex plos.template
% latex plos.template
% latex plos.template
% dvipdf plos.template
%
% % % % % % % % % % % % % % % % % % % % % %
%
% -- IMPORTANT NOTE
%
% This template contains comments intended 
% to minimize problems and delays during our production 
% process. Please follow the template instructions
% whenever possible.
%
% % % % % % % % % % % % % % % % % % % % % % % 
%
% Once your paper is accepted for publication, 
% PLEASE REMOVE ALL TRACKED CHANGES in this file and leave only
% the final text of your manuscript.
%
% There are no restrictions on package use within the LaTeX files except that 
% no packages listed in the template may be deleted.
%
% Please do not include colors or graphics in the text.
%
% Please do not create a heading level below \subsection. For 3rd level headings, use \paragraph{}.
%
% % % % % % % % % % % % % % % % % % % % % % %
%
% -- FIGURES AND TABLES
%
% Please include tables/figure captions directly after the paragraph where they are first cited in the text.
%
% DO NOT INCLUDE GRAPHICS IN YOUR MANUSCRIPT
% - Figures should be uploaded separately from your manuscript file. 
% - Figures generated using LaTeX should be extracted and removed from the PDF before submission. 
% - Figures containing multiple panels/subfigures must be combined into one image file before submission.
% See http://www.plosone.org/static/figureGuidelines for PLOS figure guidelines.
%
% Tables should be cell-based and may not contain:
% - tabs/spacing/line breaks within cells to alter layout or alignment
% - vertically-merged cells (no tabular environments within tabular environments, do not use \multirow)
% - colors, shading, or graphic objects
% See http://www.plosone.org/static/figureGuidelines#tables for table guidelines.
%
% For tables that exceed the width of the text column, use the adjustwidth environment as illustrated in the example table in text below.
%
% % % % % % % % % % % % % % % % % % % % % % % %
%
% -- EQUATIONS, MATH SYMBOLS, SUBSCRIPTS, AND SUPERSCRIPTS
%
% IMPORTANT
% Below are a few tips to help format your equations and other special characters according to our specifications. For more tips to help reduce the possibility of formatting errors during conversion, please see our LaTeX guidelines at http://www.plosone.org/static/latexGuidelines
%
% Please be sure to include all portions of an equation in the math environment.
%
% Do not include text that is not math in the math environment. For example, CO2 will be CO\textsubscript{2}.
%
% Please add line breaks to long display equations when possible in order to fit size of the column. 
%
% For inline equations, please do not include punctuation (commas, etc) within the math environment unless this is part of the equation.
%
% % % % % % % % % % % % % % % % % % % % % % % % 
%
% Please contact latex@plos.org with any questions.
%
% % % % % % % % % % % % % % % % % % % % % % % %

\documentclass[10pt,letterpaper]{article}
\usepackage[top=0.85in,left=2.75in,footskip=0.75in]{geometry}

% Use adjustwidth environment to exceed column width (see example table in text)
\usepackage{changepage}

% Use Unicode characters when possible
\usepackage[utf8]{inputenc}

% textcomp package and marvosym package for additional characters
\usepackage{textcomp,marvosym}

% fixltx2e package for \textsubscript
\usepackage{fixltx2e}

% amsmath and amssymb packages, useful for mathematical formulas and symbols
\usepackage{amsmath,amssymb}

% cite package, to clean up citations in the main text. Do not remove.
\usepackage{cite}

% Use nameref to cite supporting information files (see Supporting Information section for more info)
\usepackage{nameref,hyperref}

% line numbers
\usepackage[right]{lineno}

% ligatures disabled
\usepackage{microtype}
\DisableLigatures[f]{encoding = *, family = * }

% rotating package for sideways tables
\usepackage{rotating}

% Remove comment for double spacing
%\usepackage{setspace} 
%\doublespacing

% Text layout
\raggedright
\setlength{\parindent}{0.5cm}
\textwidth 5.25in 
\textheight 8.75in

% Bold the 'Figure #' in the caption and separate it from the title/caption with a period
% Captions will be left justified
\usepackage[aboveskip=1pt,labelfont=bf,labelsep=period,justification=raggedright,singlelinecheck=off]{caption}

% Use the PLoS provided BiBTeX style
\bibliographystyle{plos2009}

% Remove brackets from numbering in List of References
\makeatletter
\renewcommand{\@biblabel}[1]{\quad#1.}
\makeatother

% Leave date blank
\date{}

% Header and Footer with logo
\usepackage{lastpage,fancyhdr,graphicx}
\pagestyle{myheadings}
\pagestyle{fancy}
\fancyhf{}
\lhead{\includegraphics[natwidth=1.3in,natheight=0.4in]{PLOSlogo.png}}
\rfoot{\thepage/\pageref{LastPage}}
\renewcommand{\footrule}{\hrule height 2pt \vspace{2mm}}
\fancyheadoffset[L]{2.25in}
\fancyfootoffset[L]{2.25in}
\lfoot{\sf PLOS}

%% Include all macros below

\newcommand{\lorem}{{\bf LOREM}}
\newcommand{\ipsum}{{\bf IPSUM}}

%% END MACROS SECTION


\begin{document}
\vspace*{0.35in}

% Title must be 150 characters or less
\begin{flushleft}
{\Large
\textbf\newline{Article Title - Diversity dynamics in Nymphalidae butterflies: Effect of phylogenetic uncertainty on diversification rate shift estimates}
}
\newline
% Insert Author names, affiliations and corresponding author email.
\\
Carlos Pe\~na\textsuperscript{1,*},
Marianne Espeland\textsuperscript{2,\textcurrency}
\\
\bf{1} Laboratory of Genetics, Department of Biology, University of Turku, Turku, Finland
\\
\bf{2} Museum of Comparative Zoology and Department of Organismic and Evolutionary Biology, Harvard University, Cambridge, USA
\\

% Insert additional author notes using the symbols described below. Insert symbol callouts after author names as necessary.
% 
% Remove or comment out the author notes below if they aren't used.
%
% Primary Equal Contribution Note

% Current address notes
\textcurrency Current address: McGuire Center for Lepidoptera and Biodiversity,
Florida Museum of Natural History, University of Florida, Gainesville,
FL, US.
% \textcurrency b Insert current address of second author with an address update
% \textcurrency c Insert current address of third author with an address update

% Deceased author note

% Group/Consortium Author Note

* E-mail: mycalesis@gmail.com
\end{flushleft}

\subsection*{Running title: Diversity dynamics in Nymphalidae
butterflies}

% Please keep the abstract below 300 words
\section*{Abstract}
The species rich butterfly family Nymphalidae has been used to study
evolutionary interactions between plants and insects. Theories of
insect-hostplant dynamics predict accelerated diversification due to key
innovations. In evolutionary biology, analysis of maximum credibility
trees in the software MEDUSA (modelling evolutionary diversity using
stepwise AIC) is a popular method for estimation of shifts in
diversification rates. We investigated whether phylogenetic uncertainty
can produce different results by extending the method across a random
sample of trees from the posterior distribution of a Bayesian run. Using
the MultiMEDUSA approach, we found that phylogenetic uncertainty greatly
affects diversification rate estimates. Different trees produced
diversification rates ranging from high values to almost zero for the
same clade, and both significant rate increase and decrease in some
clades. Only four out of 18 significant shifts found on the maximum
clade credibility tree were consistent across most of the sampled trees.
Among these, we found accelerated diversification for Ithomiini
butterflies. We used the binary speciation and extinction model (BiSSE)
and found that a hostplant shift to Solanaceae is correlated with
increased net diversification rates in Ithomiini, congruent with the
diffuse cospeciation hypothesis. Our results show that taking
phylogenetic uncertainty into account when estimating net
diversification rate shifts is of great importance, as very different
results can be obtained when using the maximum clade credibility tree
and other trees from the posterior distribution.

\textbf{Keywords:} diversification analysis, MEDUSA, BiSSE, speciation
rate, insect-hostplant dynamics

\linenumbers

\section*{Introduction}
Hostplant shifts have been invoked to be responsible for a great part of
the biodiversity of herbivorous insects \cite{mitter1988}. The study of the
evolution of hostplant use has spawned several theories explaining the
evolutionary interactions between plants and insects \cite{nyman2012}: the
``escape-and-radiate'' hypothesis \cite{ehrlich1964}, the ``oscillation
hypothesis'' \cite{janz2011, nylin2014} or ``diffuse cospeciation'' \cite{nyman2012} and the
``musical chairs hypothesis'' \cite{hardy2014}.

The butterfly family Nymphalidae has been an important taxon for
developing some of the mentioned hypotheses. Nymphalidae contains around
6000 species \cite{van_nieukerken2011}, and several members are considered model organisms
in evolutionary biology \cite{joron2006, willmott2006, brakefield2009}.
The family most likely originated
around 94 MYA in the mid Cretaceous. Diversification of the group began
in the Late Cretaceous and most major radiations (current subfamilies)
appeared shortly after the Cretaceous-Paleogene (K-Pg) boundary
\cite{heikkila2012}. Several studies have used time-calibrated phylogenies and
diversification models to reconstruct the evolutionary history of the
group to identify patterns of accelerated or decelerated diversification
of some Nymphalidae clades \cite{heikkila2012, elias2009, fordyce2010, wahlberg2009}.
For example, it has been
suggested that climate change in the Oligocene and the subsequent
diversification of grasses has led to diversification of the subfamily
Satyrinae \cite{pena2008} due to the abundance of grasses over extensive
geographic areas (``resource abundance-dependent diversity dynamics''
hypothesis): Fordyce (2010) \cite{fordyce2010} found increased net diversification
rates in some Nymphalidae lineages after a major hostplant shift, which
appears to be in agreement with the ``escape-and-radiate'' model of
diversification \cite{ehrlich1964}.

Although it has been suggested that part of the great diversity of
Nymphalidae butterflies is a result of hostplant-insect dynamics, it is
necessary to use modern techniques to investigate whether the
diversification patterns of Nymphalidae are in agreement with the
theoretical predictions. It is necessary to test whether the overall
diversification pattern of Nymphalidae is congruent with events of
sudden diversification bursts due to hostplant shift or climatic events
\cite{nylin2014, ferrer2013}.

In this study, we used a time-calibrated genus-level phylogenetic
hypothesis for Nymphalidae butterflies \cite{wahlberg2009} to investigate patterns
of diversification. We applied the statistical method MEDUSA
\cite{alfaro2009, harmon2011}, to study the diversification pattern of Nymphalidae
butterflies. MEDUSA fits likelihood models of diversification into a
time-calibrated tree and tests whether allowing increases or decreases
in speciation and extinction rates within the tree produces better fit
of the models. MEDUSA is able to take into account unsampled extant
species diversity during model fitting and it is normally applied to the
maximum clade credibility phylogenetic tree. Particularly, we wanted to
study the effects of phylogenetic uncertainty and by using the extended
MEDUSA method called MultiMEDUSA \cite{alfaro2009}. We also tested whether
hostplant association dynamics can explain the diversification patterns
of component Nymphalidae lineages by testing whether character states of
hostplant use affected the diversification pattern of those lineages
employing the binary speciation and extinction model (BiSSE) as
implemented in the R package \texttt{diversitree} \cite{fitzjohn2012}.


% You may title this section "Methods" or "Models". 
% "Models" is not a valid title for PLoS ONE authors. However, PLoS ONE
% authors may use "Analysis" 
\section*{Materials and Methods}
\subsection*{Data}

For analyses, we used the phylogenetic trees from the study of Wahlberg
et al. (2009) \cite{wahlberg2009} that were generated using DNA sequence data from
10 gene regions for 398 of the 540 valid genera in Nymphalidae. We
employed the maximum clade credibility tree (MCC tree) (Fig. 1) as well
as a random sample of 1000 trees from their BEAST run after burnin
\cite{wahlberg2009}. Their original BEAST run was for 40 million generations. We
used a burnin of 25 million generations and took a random sample of 1000
trees using Burntrees v.0.1.9 \url{http://www.abc.se/~nylander/} (S6 File)
in order to account for phylogenetic uncertainty when performing
the diversification analyses.

\begin{figure}[h]
\caption{{\bf Results of the MEDUSA analysis run on the maximum
clade credibility tree from Wahlberg et al. (2009)}. Rate shifts were estimated for the
following nodes (besides the background rate): 1) root, 2) Limenitidinae +
Heliconiinae, 3) \emph{Ypthima}, 4) \emph{Charaxes}, 5) Ithomiini in
part, 6) Satyrinae, 7) Coenonymphina 8) Phyciodina in part, 9) Danaini in
part, 10) \emph{Coenonympha}, 11) \emph{Caeruleuptychia} +
\emph{Magneuptychia}, 12) \emph{Callicore} + \emph{Diaethria}, 13)
Satyrina, 14) Mycalesina, 15) \emph{Pedaliodes}, 16) \emph{Dryas} +
\emph{Dryadula}, 17) \emph{Taenaris}, 18) Pseudergolinae, 19)
\emph{Anaeomorpha} + \emph{Hypna}. Circles on
nodes indicate the diversification shift number as found by MEDUSA.
Numbers next to circles indicate the posterior probability values for
such nodes.}
\label{fig1}
\end{figure}

Species richness data for Nymphalidae genera were compiled from several
sources including the specialist-curated lists 
\cite{tolweb2007, lamas2004} and curated lists of the
Global Butterfly Names project \cite{gbnp}.
We assigned the species numbers of genera not included in the phylogeny
to the closest related genus that was included in Wahlberg et al. (2009)
\cite{wahlberg2009} study according to available phylogenetic studies 
\cite{matos2013, brower2010, kodandaramaiah2010, kodandaramaiah2010a,
ortiz2013, desilva2010, freitas2004, pena2006, penz1999, silva2008,
pena2011, pena2010},
taxonomical classification and morphological resemblance when no
phylogenies were available (S7 Dataset).

Hostplant data for Nymphalidae species were compiled from several
sources including \cite{ackery1988}, HOSTS database (\url{http://bit.ly/YI7nwW}),
\cite{dyer2002, beccaloni2008, janzen2009} and others (S8--S9 Datasets) for a
total of 6586
hostplant records, including 428 Nymphalidae genera and 143 plant
families and 1070 plant genera. It was not possible to find any
hostplant data for 35 butterfly genera.



\subsection*{Analyses of Diversification}

We used the statistical software R version 3.0.1 \cite{r2013} in combination
with the APE \cite{popescu2012}, GEIGER \cite{harmon2008}, MEDUSA
\cite{alfaro2009} and
\texttt{diversitree} \cite{fitzjohn2012} packages along with our own scripts to
perform the analyses (included as S10--S12 Supporting Information). All
analyses were run on the 1000 random trees from Wahlberg et al. (2009)
\cite{wahlberg2009}, on the maximum clade credibility tree derived from these and
the MCC tree \cite{wahlberg2009}.


\subsection*{Detecting diversification shifts on phylogenetic trees}

Patterns of diversification in Nymphalidae were analyzed by using MEDUSA
version 093.4.33 \cite{alfaro2009} on the maximum clade credibility tree from
Wahlberg et al. (2009) \cite{wahlberg2009}. MEDUSA fits alternative birth-death
likelihood models to a phylogenetic tree in order to estimate changes in
net diversification rates along branches. MEDUSA estimates likelihood
and AIC scores for the simplest birth-death model, with two parameters,
the rates \texttt{r}: net diversification \texttt{=} speciation
(\texttt{b}) \texttt{-} extinction (\texttt{d}) and (\(\varepsilon\):
relative extinction \texttt{= d/b}.
The maximum-likelihood values are recalculated when each branch in the tree
is selected to contain an additional breakpoint (additional parameters
\texttt{b} and \texttt{d}) in net diversification rate. 
The process is repeated by fitting additional breaking points to the tree.
The fit of  alternative models is compared using AIC
scores and stop when the improvement of the incrementally more complex models 
is smaller than a cutoff value which
depends of the number of terminals.
MEDUSA finds the likelihood of the models after taking into
account branch lengths and number of species per lineage \cite{alfaro2009}. Most
studies using MEDUSA only run the method on a single tree, usually the
maximum clade credibility tree, which makes the assumption that this
tree is correct. We wanted to study the effects of phylogenetic
uncertainty on estimation of net diversification rate shifts and
therefore selected 1000 random genus-level trees from the posterior
distribution of the Bayesian run from Wahlberg et al., 2009 \cite{wahlberg2009}
(MultiMEDUSA, S6 File). We calculated a new MCC tree from the
selection of trees, ran MEDUSA and MultiMEDUSA on the selection of 1000
trees and summarized the estimated changes in net diversification rates
for nodes across all trees. Patterns of change in net diversification
rates are considered significant if they are found at the same node in
at least 90\% of the trees. We also expected to find similar values of
net diversification (\texttt{r = b - d}) and relative extinction rate
(\(\varepsilon\) \texttt{= d/b}) values across the 1000 trees for the
nodes where changes in diversification tempo occurs. We used the AICc
threshold of 7.8 units, as estimated by MEDUSA, as the limit for a
significantly better fit to select among increasingly complex
alternative models.

\subsection*{Estimation of trait-dependent speciation rates}

As the MEDUSA and MultiMEDUSA approaches estimated an increase in net
diversification in the clade Ithomiini, we tested whether this pattern
can be explained by an increase in the birth-rate due to hostplant use
and performed analyses using the BiSSE model \cite{maddison2007} 
as implemented in the R package
\texttt{diversitree} \cite{fitzjohn2012}. The Markov Chain Monte Carlo algorithm
was run for 10000 generations discarding the first 7500 as burnin.
Because most of Nymphalidae butterflies are restricted to
use one plant family as hostplant (232 genera use only one plant family
versus 176 that use more than one; based on our data in S8 Dataset),
the character states can be coded as presence/absence, for which the
BiSSE analysis is better suited. BiSSE was designed to test whether a
binary character state has had any effect on increased net
diversification rate for a clade \cite{maddison2007}. We used our compiled data
of hostplant use to produce a binary dataset for the character ``feeding on
the plant family Solanaceae'' with two states (absence = 0; presence =
1) (S13 Dataset). The family Solanaceae include the hostplants of the
diverse Ithomiini butterflies and closest relatives \cite{willmott2006}. 
We analyzed the data using BiSSE employing the Markov Chain
Monte Carlo algorithm on the maximum clade credibility tree, taking into
account missing taxa by using the parameter ``sampling factor''
(\texttt{sampling.f}) in \texttt{diversitree}. 
We also run a BiSSE analysis
to test whether the trait ``feeding on Apocynaceae'' had any effect on
net diversification rates found for Nymphalidae lineages feeding on
Apocynaceae.
We also used likelihood ratio tests to evaluate the influence of
Solanaceae-feeding on diversification rates by comparing a constrained versus
a relaxed model of diversification.
The constrained model assumes no effect of hostplant use on diversification by
enforcing equal speciation rates across both character states. In the relaxed
model, all speciation and extinction parameters are estimated for each
character state.
The analyses were run across a sample of 250
trees from the selected 1000 random trees from the posterior
distribution. The records of \emph{Vanessa} and \emph{Hypanartia}
feeding on Solanaceae \cite{beccaloni2008, scott1986} might be incorrect as it is unlikely
that these species are able to feed on members of this plant family. We tested
whether coding these two genera as not feeding on Solanaceae affected
our results. We expected that coding these two genera as either absence
or presence would not have a significant effect on our results. It has
been shown that BiSSE performs poorly under certain conditions \cite{davis2013}.
However, our data has adequate number of taxa under analysis (more than
300 tips), adequate speciation bias (between 1.5x and 2.0x), character
state bias (around 8x) and extinction bias (around 4x) for the analysis
of Solanaceae hostplants. Thus, BiSSE is expected to produce robust
results \cite{davis2013}.

\section*{Results}

\subsection*{Detecting diversification shifts on the maximum clade
credibility tree}

The MEDUSA analysis on the MCC tree in combination with richness data
estimated 18 significant changes in the tempo of diversification in
Nymphalidae history (Fig. 1; Table 1). The corrected AICc acceptance
threshold for adding subsequent piecewise birth-death processes to the
overall model was set to 7.8 units, as prescribed by MEDUSA. In all
MEDUSA analyses, the maximum number of inferred diversification shifts
in all trees was 26. The background net diversification rate for
Nymphalidae was estimated as \texttt{r = 0.092} lineages per Million of
years and the AICc score for the best fit model was \texttt{5090.5}
(Table 1). MEDUSA also estimated that the basic constant birth-death
model was not a better explanation for our data (AICc
\texttt{= 5449.3}).

% Table 1
\begin{table}[!ht]
\begin{adjustwidth}{-2.25in}{0in}
\caption{{\bf Significant net diversification rate shifts found in the MEDUSA analysis of the Nymphalid maximum clade credibility tree.}}
\begin{tabular}{|l|c|c|r|r|c|l|}
\hline
Shift N$^\circ$ & Shift.Node & Model & r          & LnLik.part & AICc     & Taxa                                                       \\ \hline
1               & 399        & yule  & 0.092459   & -1055.957  &          & Nymphalidae (root)                                         \\ \hline
2               & 691        & bd    & 0.054129   & -406.3703  &          & Limenitidinae + Heliconiinae                               \\ \hline
3               & 299        & yule  & 0.311199   & -6.3058    &          & \emph{Ypthima}                                             \\ \hline
4               & 224        & yule  & 0.290989   & -6.2601    &          & \emph{Charaxes}                                            \\ \hline
5               & 750        & yule  & 0.186913   & -147.4146  &          & Oleriina + Ithomiina + Napeogenina + Dircennina + Godyrina \\ \hline
6               & 405        & yule  & 0.116252   & -555.0276  &          & Satyrinae                                                  \\ \hline
7               & 495        & yule  & 0.064656   & -124.9143  &          & Coenonymphina                                              \\ \hline
8               & 609        & yule  & 0.240562   & -35.0908   &          & Phyciodina in part                                         \\ \hline
9               & 787        & bd    & 0.042332   & -43.7819   &          & Danaini in part                                            \\ \hline
10              & 231        & yule  & 0.209416   & -4.7955    &          & \emph{Coenonympha}                                         \\ \hline
11              & 478        & yule  & 0.311684   & -9.2218    &          & \emph{Caeruleuptychia} + \emph{Magneuptychia}              \\ \hline
12              & 659        & yule  & 0.219253   & -11.2686   &          & \emph{Callicore} + \emph{Diaethria}                        \\ \hline
13              & 444        & yule  & 0.220615   & -43.7812   &          & Satyrina                                                   \\ \hline
14              & 524        & yule  & 0.190754   & -26.0651   &          & Mycalesina                                                 \\ \hline
15              & 355        & yule  & 0.234041   & -6.2013    &          & \emph{Pedaliodes}                                          \\ \hline
16              & 714        & yule  & 0          & 0          &          & \emph{Dryas} + \emph{Dryadula}                             \\ \hline
17              & 377        & yule  & 0.311671   & -4.1986    &          & \emph{Taenaris}                                            \\ \hline
18              & 688        & yule  & 0.024724   & -17.5256   &          & Pseudergolinae                                             \\ \hline
19              & 583        & yule  & 0          & 0          & 5090.492 & \emph{Anaeomorpha} + \emph{Hypna}                          \\ \hline
\end{tabular}
\begin{flushleft}Shift.Node=node number, Model=preferred diversification model by MEDUSA, r=net diversification rate, LnLik.part=log likelihood value.
\end{flushleft}
\label{table1}
\end{adjustwidth}
\end{table}

Some of the 18 changes in diversification correspond to rate increases
in very species-rich genera: \emph{Ypthima} (\texttt{r = 0.311}),
\emph{Charaxes} (\texttt{r = 0.291}), \emph{Callicore} +
\emph{Diaethria} (\texttt{r = 0.220}), \emph{Pedaliodes}
(\texttt{r = 0.196}) and \emph{Taenaris} (\texttt{r = 0.234}). We found
rate increases for other clades as well such as: Mycalesina
(\texttt{r = 0.191}), Oleriina + Ithomiina + Napeogenina + Dircennina +
Godyridina (\texttt{r = 0.187}), Satyrinae (\texttt{r = 0.116}),
Phyciodina in part (\texttt{r = 0.241}) and Satyrina
(\texttt{r = 0.221}), \emph{Coenonympha} (\texttt{r = 0.209}),
\emph{Caeruleuptychia} + \emph{Magneuptychia} (\texttt{r = 0.312}) and
\emph{Taenaris} (\texttt{r = 0.312}). We also found decreases in net
diversification rates for Limenitidinae + Heliconiinae
(\texttt{r = 0.0541}), part of Danaini (\texttt{r = 0.0423}),
Pseudergolinae (\texttt{r = 0.024}) and Coenonymphina
(\texttt{r = 0.065}) (Table 1).


\subsection*{Phylogenetic uncertainty in the MultiMEDUSA
approach}

We used MEDUSA to find out whether taking into account the phylogenetic
signal from the random sample of 1000 trees from the posterior
distribution can return similar estimates of diversification to the
values obtained from the MCC tree. We ran MultiMEDUSA on the random
sample of 1000 trees (S6 File) from the posterior distribution and
compared the results with a MEDUSA analysis on the MCC tree derived from
this sample (S14 Supporting Information).

We found that the analysis by MultiMEDUSA on the 1000 trees estimated
lower median net diversification rates for the diversification shifts
found by MEDUSA on the MCC tree derived from the random sample of trees
(Table 2). Although the diversification pattern found by MEDUSA and
MultiMEDUSA was the same, the latter consistently estimated lower rates.
Furthermore, the shifts recovered with low net diversification rate on
the MCC were recovered with negative net diversification rate by
MultiMEDUSA. The background diversification and all shifts found by
MEDUSA on the 1000 trees are provided as an R object in S15 File.

% Table 2
\begin{table}[!ht]
\begin{adjustwidth}{-2.25in}{0in} % Comment out/remove adjustwidth environment if table fits in text column.
\caption{{\bf Differences in rates estimated by MEDUSA on the MCC
tree from the sample of trees from the posterior distribution and the
MultiMEDUSA approach. Shift consistently recovered across the
sample of trees in bold face.}}
\begin{tabular}{|l|c|r|r|}
    \hline
Shift.Node & rate by MEDUSA & Median rate by MultiMEDUSA & probability of being recovered \\ \hline
1          & 0.092          & not found                  & 0.000                          \\ \hline
2          & 0.055          & -0.030                     & 0.864                          \\ \hline
\bf{3}     & 0.184          & \bf{0.219}                 & \bf{0.961}                     \\ \hline
\bf{4}     & 0.166          & \bf{0.212}                 & \bf{0.996}                     \\ \hline
\bf{5}     & 0.111          & \bf{0.101}                 & \bf{0.927}                     \\ \hline
6          & 0.119          & 0.039                      & 0.131                          \\ \hline
7          & 0.066          & -0.052                     & 0.195                          \\ \hline
8          & 0.232          & 0.166                      & 0.319                          \\ \hline
9          & 0.042          & -0.049                     & 0.897                          \\ \hline
10         & 0.058          & 0.149                      & 0.619                          \\ \hline
11         & 0.311          & 0.208                      & 0.831                          \\ \hline
\bf{12}    & 0.219          & \bf{0.135}                 & \bf{0.911}                     \\ \hline
13         & 0.082          & 0.117                      & 0.276                          \\ \hline
14         & 0.099          & 0.115                      & 0.379                          \\ \hline
15         & 0.113          & 0.127                      & 0.825                          \\ \hline
16         & 0.222          & -0.005                     & 0.659                          \\ \hline
17         & 0.243          & 0.248                      & 0.785                          \\ \hline
18         & 0.192          & -0.064                     & 0.024                          \\ \hline
19         & 0.064          & -0.087                     & 0.381                          \\ \hline
\end{tabular}
\begin{flushleft}
\end{flushleft}
\label{table2}
\end{adjustwidth}
\end{table}

We also compared the results from MultiMEDUSA (derived from the sample
of 1000 trees) with the shifts found by MEDUSA on the MCC tree derived
from this random sample. In the summary statistics, MultiMEDUSA reports
the frequency of the diversification shifts found in the trees
(parameter \texttt{sum.prop}). Thus, if a node is found in only half of
the 1000 trees, but the phylogenetic signal was strong enough to be
picked up by MEDUSA and a node shift was found most of the time, then
the \texttt{sum.prop} should be close to 1. For example, the
\texttt{Charaxes} + \texttt{Polyura} clade was found in only 256 trees,
however MultiMEDUSA was consistently able to find a diversification
shift for that node and the \texttt{sum.prop} value is 0.996.

For the diversification shifts found in both the MCC tree and most of
the samples of 1000 trees (frequency more than 90\%; Table 2), the
MultiMEDUSA approach recovered different rates of diversification than
those found when the MCC tree alone was used.

There were four net diversification rate increases found in the trees from the
random sample (Table 2): (i) the
genus \emph{Ypthima} (\(r = 0.22\)); (ii) the genus \emph{Charaxes}
(\(r = 0.21\)); (iii) Ithomiini subtribes Oleriina + Ithomiina + Napeogenina +
Dircennina + Godyridina (\(r = 0.10\)); and (iv) the clade of 
\emph{Callicore} + \emph{Diaethria} (\(r = 0.135\)).

MultiMEDUSA provided mean and standard deviation statistics for the
diversification values on the shifts on the 1000 trees 
(S16--S17 Text and Figure), and found that the estimated net diversification rate
values had great variation across the posterior distribution of trees. A
boxplot of the net diversification rate values estimated for the clades
that appear in the MCC tree shows that some shifts are estimated as
increased or slowed diversification pace depending on the tree used for
analysis (Fig. 2). This variation is especially wide for the clade
formed by the genera \emph{Magneuptychia} and \emph{Caeruleuptychia}
because MEDUSA estimated diversification values from six times the
background net diversification rate (\texttt{r = 0.5234}) to almost zero
(\texttt{r = 1.22e-01}). The rates for \emph{Taenaris} were between 0.14
and 0.44 (mean value 0.25). Similar degrees of variation were found in
the nodes for \emph{Ypthima}, \emph{Charaxes} and \emph{Coenonympha}
(Fig. 2). The net diversification rates estimates for the clades
(Oleriina + Ithomiina + Napeogenina + Dircennina + Godyridina),
Limenitidinae + Heliconiinae and Pseudergolinae are relatively
consistent across the 1000 trees (Fig. 2).

\begin{figure}[h]
\caption{{\bf Diversification rates for taxa estimated by MultiMEDUSA on
the samples of 1000 random trees}}
\label{fig2}
\end{figure}

It is also evident that not all the diversification shifts estimated on
the MCC tree are consistently recovered in most of the 1000 trees. Some
of the shifts in the MCC tree are recovered in very few trees, for
example the shift for the clade Satyrinae is recovered with a
probability of 0.18 (Fig. 3).

\begin{figure}[h]
\caption{{\bf Results of MultiMEDUSA analysis showing the
probability of specific nodes being characterized by significant shifts
in diversification rate}}
\label{fig3}
\end{figure}

\subsection*{Estimation of trait-dependent speciation
rates}

The MEDUSA analyses, taking into account phylogenetic uncertainty,
estimated a net diversification rate increase in part of the clade
Ithomiini across more than 95\% of the trees. Our BiSSE analysis found a
positive effect of the character state ``feeding on Solanaceae'' on the
net net diversification rate on part of Ithomiini (Oleriina + Ithomiina
+ Napeogenina + Dircennina + Godyridina) (Fig. 4). The estimated mean
net diversification rate for taxa that do not feed on Solanaceae was
\texttt{r = 0.11} while the net diversification rate for the Solanaceae
feeders was \texttt{r = 0.16} (see S1 Fig. for a boxplot of speciation
and extinction values for the 95\% credibility intervals). The same
analysis considering \emph{Vanessa} and \emph{Hypanartia} as
non-Solanaceae feeders due to dubious records produced the same pattern
and net diversification rates (S2 Fig.). Therefore, the rest of the
analyses were performed assuming these two genera as Solanaceae feeders.
We constrained the BiSSE likelihood model to force equal rates of
speciation for both character states in order to test whether the model
of different speciation rates is a significantly better explanation for
the data. A likelihood ratio test found that the model for increased net
diversification rate for nymphalids feeding on Solanaceae is a
significantly better explanation than this character state having no
effect on diversification (\(\chi^2 = 12.3; 1 df; p < 0.001\)) (Table 3;
character states available in S13 Dataset, code in S18 Supporting Information,
and mcmc run in S19 Supporting Information). We combined the post-burnin mcmc generations
from running BiSSE on 250 trees from the random sample of 1000 trees and
found the same pattern as the BiSSE analysis on the maximum clade
credibility tree (combined mcmc run in S20 Supporting Information; profiles
plot of speciation rates in S3 Fig.; boxplot of 95\% credibility intervals in
S4 Fig.). A BiSSE analysis to test whether the trait ``feeding on
Apocynaceae'' had any effect on increased net diversification rates
found similar speciation rates for lineages feeding on Apocynaceae and
other plants (S5 Fig.).

% Table 3
\begin{table}[!ht]
% \begin{adjustwidth}{-2.25in}{0in} % Comment out/remove adjustwidth environment if table fits in text column.
    \caption{{\bf Likelihood ratio test between the model of increased
diversification of nymphalids feeding on Solanaceae against a model
forcing equal speciation rates (no effect on diversification).}}
\begin{tabular}{|l|c|r|r|l|r|}
    \hline
Df           & lnLik & AIC     & ChiSq  & \multicolumn{1}{r}{p} &         \\ \hline
full         & 6     & -1613.3 & 3238.5 &                       &         \\ \hline
equal.lambda & 5     & -1619.4 & 3248.9 & 12.3                  & 0.00045 \\ \hline
\end{tabular}
\begin{flushleft}
\end{flushleft}
\label{table3}
% \end{adjustwidth}
\end{table}

\begin{figure}[h]
\caption{{\bf BiSSE analysis of diversification of nymphalids due
to feeding on Solanaceae hostplants}.
Speciation and net diversification
rates are significantly higher in Solanaceae feeders (speciation rate =
$\lambda1$, net diversification rate = $r1$}
\label{fig4}
\end{figure}

\section*{Discussion}

\subsection*{Effects of phylogenetic uncertainty on the performance of
MEDUSA}

The MEDUSA method has been used to infer changes in net diversification
rates in a phylogenetic tree. Since its publication \cite{alfaro2009} the results
of using MEDUSA on a single tree, the maximum clade credibility tree,
have been used for generation of hypotheses and discussion
\cite{heikkila2012, litman2011, ryberg2012}. However, different diversification
shifts and different
rates of diversification are found for certain lineages when
phylogenetic uncertainty was taken into account by using MEDUSA on a
random sample of trees from the posterior distribution of a Bayesian
run. We found that some diversification shifts, estimated on the
Nymphalidae maximum clade credibility tree, were found with very low
probability in the random sample of 1000 trees from the posterior
distribution (Fig. 3, Table 2). We also found that, even though the
analyses estimated the same diversification shifts on two or more trees,
the estimated net diversification rates could vary widely (Fig. 2). For
example, in our Nymphalidae trees, we found that the shift for
\emph{Magneuptychia} and \emph{Caeruleuptychia} had a rate from
\texttt{r = 0.5234}, higher than the background net diversification
rate, to almost zero. This means that observed patterns and conclusions
can be completely contradictory depending on tree choice.

In this study, the effect of phylogenetic uncertainty on the inferred
diversification shifts by MEDUSA is amplified because some Nymphalidae
taxa appear to be strongly affected by long-branch attraction artifacts
\cite{pena2011}. Thus, the Bayesian runs are expected to recover alternative
topologies on the posterior distribution of trees, resulting in low
support and posterior probability values for the nodes. For example,
posterior probability values for clades in Satyrinae are very low
\cite{wahlberg2009}. As a result, MEDUSA inferred a net diversification rate
increase for the Satyrinae (which includes Satyrini and its sister
clade) in the maximum clade credibility tree, but this was recovered in
the MultiMEDUSA analysis in only 13\% of the random sample of trees.

If there is strong phylogenetic signal for increases or decreases in net
diversification rates for a node, it is expected that these shifts would
be inferred by MEDUSA in most of the posterior distribution of trees.
However, weak phylogenetic signal for some nodes can cause some clades
to be absent in some trees and MEDUSA will be unable to estimate any
diversification shift (due to a non-existent node).
The relatively low phylogenetic support for many nodes in the Nymphalidae
tree is likely the reason
why MEDUSA estimated net diversification rate shifts with a probability
higher than 0.90 in the sample of trees for only four shifts: the genus
\emph{Charaxes}, the genus \emph{Ypthima}, part of Ithomiini and the
clade \emph{Callicore} + \emph{Diaethria} (Fig. 3), while estimating
shifts for other lineages with much lower probability.


\subsection*{Hostplant use and diversification in Nymphalidae}

\paragraph*{Ithomiini}

Keith Brown suggested that feeding on Solanaceae was an important event
in the diversification of Ithomiini butterflies \cite{brown1987}. Ithomiini
butterflies are exclusively Neotropical and most species feed on
Solanaceae hostplants during larval stage \cite{willmott2006}. Optimizations of the
evolution of hostplant use on phylogenies evidence a probable shift from
Apocynaceae to Solanaceae in the ancestor of the tribe \cite{willmott2006, brower2006}.
Fordyce (2010) \cite{fordyce2010} found that the Gamma statistics, a LTT plot of
an Ithomiini phylogeny and the fit of the density-dependent model of
diversification are consistent with a burst of diversification in
Ithomiini following the shift from Apocynaceae to Solanaceae.

We investigated whether the strong signal for an increase in net
diversification rate for Ithomiini (found by MEDUSA) can be explained
due to the use of Solanaceae plants as hosts during larval stage. For
this, we used a Bayesian approach \cite{fitzjohn2009} to test whether the trait
``feeding on Solanaceae'' had any effect on the diversification of the
group.

Our BiSSE analysis, extended to take into account missing taxa and
phylogenetic uncertainty, shows a significantly higher net
diversification rate for Ithomiini taxa, which can be attributed to the
trait ``feeding on Solanaceae hostplants'' (Fig. 4). This is in
agreement with previous findings using other statistical methods
\cite{fordyce2010}. Due to the fact that Ithomiini are virtually the only
nymphalids using Solanaceae as hostplants, it is possible that the trait
responsible for a higher diversification of Ithomiini might not be the
hostplant character. As noted by Maddison et al. (2007) \cite{maddison2007}, the
responsible trait might be a codistributed character such as a trait
related to the ability to digest secondary metabolites.

Solanaceae plants contain chemical compounds and it has been suggested
that the high diversity of Ithomiini is consistent with the
``escape-and-radiate scenario'' due to a shift onto Solanaceae
\cite{fordyce2010}
and radiation scenarios among chemically different lineages of
Solanaceae plants \cite{willmott2006, brown1987}. According to this theory, the shift from
Apocynaceae to Solanaceae allowed Ithomiini to invade newly available
resources due to a possible key innovation that allowed them cope with
secondary metabolites of the new hosts. Additional studies are needed to
identify the actual enzymes that Ithomiini species might be using for
detoxification of ingested food as they have been found in other
butterfly groups \cite{wheat2007}.

The increase in diversification rate inferred by MEDUSA occurred after the
probable shift from Apocynaceae to Solanaceae, as the Solanaceae feeders
in the subtribes Melinaeina and Mechanitina are not included in the
diversification shift (shift number 5 in Fig. 1). The apparent
conflicting results from MEDUSA and BiSSE can be explained by the low
species-richness of the subtribes Melinaeina and Mechanitina compared to
the other subtribes included in the shift (52 versus 272 species). It
can be that MEDUSA is more conservative than BiSSE and is not including
Melinaeina and Mechanitina in the shift due to low species numbers.

Although the Solanaceae genera used by the Ithomiini clades are well
known \cite{willmott2006}, we do not have any understanding on the physiological
routes involved in the detoxification of Solanaceae compounds by the
several lineages of Ithomiini. We can speculate that older lineages
exploiting a novel toxic resource \cite{willmott2006, wahlberg2009} 
may be inefficient
in metabolizing plant toxins and that younger lineages are able to deal
with toxins more efficiently, so that host switching events within
Solanaceae are possible, which can lead to higher diversification.
Studies in \emph{Papilio} species have reported that detoxification
enzymes can become more efficient in metabolizing toxins than ancestral
configurations of the proteins, providing more opportunities for
hostplant switches \cite{li2003}. This might be the reason why the basal
Ithomiini subtribes Melinaeina and Mechanitina are so species-poor and
restricted to few Solanaceae hosts \cite{willmott2006}, while recent subtribes are
species-rich and have expanded their host range into several Solanaceae
lineages \cite{willmott2006}. It might be that the switch to feeding in Solanaceae
was an important event in the evolutionary history of Ithomiini, but the
actual radiation occurred after critical physiological changes (a
probable key innovation) allowed efficient detoxification of Solanaceae
toxins.

The diffuse cospeciation hypothesis predicts almost identical ages of
insects and their hostplants, while the ``resource abundance-dependent
diversity'' and the ``escape-and-radiate'' hypotheses posit that insects
diversify after their hostplants \cite{nyman2012, ehrlich1964, janz2011}.
Wheat et al. (2007)
\cite{wheat2007} found strong evidence for a model of speciation congruent with
Ehrlich and Raven's hypothesis in Pieridae butterflies due to, in
addition to the identification of a key innovation, a burst of
diversification in glucosinolate-feeding taxa shortly afterwards (with a
lag of circa 10 MY). According to a recent dated phylogeny of the
Angiosperms \cite{bell2010}, the family Solanaceae split from its sister group
about 59 (49--68) MYA and diversification started (crown group age)
around 37 (29--47) MYA. Wahlberg et al. (2009) \cite{wahlberg2009} give the ages
for origin and diversification for Ithomiini at 45 (39--53) and 37
(32--43) MYA, respectively. Thus, current evidence shows that Solanaceae
and Ithomiini might have diversified around the same time, during the
Late Eocene and Oligocene, and this would be congruent with the diffuse
cospeciation hypothesis.

\paragraph*{Danaini}
Our MultiMEDUSA approach showed a significant slowdown in net
diversification rate in the subtribe Danaina of the Danini. Both Danaina
and the sister clade Euploeina feed mainly on Apocynaceae and thus a
hostplant shift should not be responsible for the observed slowdown of
diversification in the Danaina. As expected, our BiSSE analysis of
Apocynaceae feeders shows that there is no effect of feeding on this
plant family on the net diversification rates of Nymphalidae lineages.
Many of the Danaina are large, strong fliers, highly migratory and
involved in mimicry rings, including the best-known migratory butterfly,
the monarch (\emph{Danaus plexippus}).
The causes for a lower net diversification rate in the
Danaina remains to be investigated, but their great dispersal power
might be involved in preventing allopatric speciation. It has been found
in highly vagile species in the nymphalid genus \emph{Vanessa} that
dispersal has homogenized populations due to gene flow, as 
vagile species seem to be genetically homogeneous among
populations \cite{wahlberg2011}.

\paragraph*{Satyrinae}

Lineages in the diverse family Satyrinae radiated simultaneously with
the radiation of their main hostplant, grasses, during the climatic
cooling in the Oligocene \cite{pena2008}. Thus, it is somewhat surprising that
part of Satyrinae were found to have accelerated diversification in only
13\% of the trees from the posterior distribution. Although this can be
attributed to low phylogenetic signal \cite{wahlberg2009}, the clade Satyrini is
very robust \cite{wahlberg2009} and MEDUSA failed to identify any significant
accelerated net diversification rate for Satyrini.
It appears that the
radiation of Satyrini as a whole was not remarkably fast and therefore
not detected by MEDUSA, although it estimated a diversification shift
for Satyrini + its sister clade. This should be expected if the
diversification of Satyrini occurred in a stepwise manner, with pulses
or bursts of diversification for certain lineages but unlikely for the
tribe Satyrini as a whole.

\section*{Conclusions}
We found that even though MEDUSA estimated several diversification
shifts in the maximum clade credibility tree of Nymphalidae, only a few
of these shifts were found in more than 90\% of the trees from the
posterior distribution. In the literature, it is common practice that
conclusions are based on the shifts estimated on the maximum clade
credibility tree. However, by using a MultiMEDUSA approach, we found
that for this Nymphalidae dataset some of these shifts might be greatly
affected by phylogenetic uncertainty. Moreover, some of these shifts can
be recovered either as increases or decreases in net diversification
rate depending on the tree from the posterior distribution that was used
for analysis. This means that contradictory conclusions would be made if
only the maximum clade credibility tree was used for analysis.
We recommend that all datasets should be analyzed using the MultiMEDUSA
approach in order to test whether the results are robust
when phylogenetic uncertainty is taken into account.

MEDUSA appears to be sensitive to the number of nodes with high
posterior probability and width of age confidence intervals. For our
data, it would be necessary to obtain a posterior distribution of trees
with no conflicting topology, and very similar estimated ages for nodes
in order to consistently recover most of the diversification shifts on
the posterior distribution of trees that were inferred by MEDUSA on the
MCC tree.

Our MultiMEDUSA approach to perform analyses on the posterior
distribution of trees found strong support for an increase in net
diversification rate in the tribe Ithomiini, the genus \emph{Ypthima},
genus \emph{Charaxes}, the clade \emph{Callicore} + \emph{Diaethria} and for a
decrease in net diversification rate in the
subtribe Danaina. Due to phylogenetic uncertainty, we did not obtain
strong support for other diversification shifts in Nymphalidae. Our
BiSSE analysis corroborated other studies in that Solanaceae-feeding,
or a codistributed character, was likely important in the
diversification of Ithomiini butterflies. However, by applying MEDUSA we
found that a critical character in the radiation of the group might have
appeared after the shift from Apocynaceae to Solanaceae. We also found
that Apocynaceae-feeding is not responsible for the
slowdown of diversification in Danaina. Ithomiini and Solanaceae
diversified near simultaneously, which is in agreement with the diffuse
cospeciation hypothesis \cite{nyman2012, janz2011}.

\section*{Supporting Information Legends}
\subsection*{S1 Fig.}
\label{S1_Fig.}
{\bf Boxplot of speciation (\(\lambda\)) and extinction
(\(\mu\)) values for the 95\% credibility intervals of values estimated
by BiSSE analysis of diversification due to feeding on Solanaceae
plants}.

\subsection*{S2 Fig.}
\label{S2_Fig.}
{\bf BiSSE analysis of diversification of nymphalids
due to feeding on Solanaceae hostplants assuming \emph{Vanessa} and
\emph{Hypanartia} as non-Solanaceae feeders. The same pattern is
recovered, speciation and net diversification rates are significantly
higher for Solanaceae feeders \(\lambda1\), r1)}.

\subsection*{S3 Fig.}
\label{S3_Fig.}
{\bf Net diversification rates of nymphalids feeding on
Solanaceae plants as estimated by combining post-burnin runs of BiSSE on
the 1000 trees from the posterior distribution}.

\subsection*{S4 Fig.}
\label{S4_Fig.}
{\bf Boxplot of speciation and extinction values for
the 95\% credibility intervals of values estimated by BiSSE analysis of
diversification due to feeding on Solanaceae plants on the combined
post-burnin runs on 1000 trees from the posterior distribution.}.

\subsection*{S5 Fig.}
\label{S5_Fig.}
{\bf BiSSE analysis of diversification of nymphalids
due to feeding on Apocynaceae hostplants}. Speciation and net
diversification rates are similar.

\subsection*{S6 File}
\label{S6_File}
{\bf 1000 random trees from \cite{wahlberg2009}}.

\subsection*{S7 Dataset}
\label{S7_Dataset}
{\bf Species richness for lineages in Nymphalidae}.

\subsection*{S8Dataset}
\label{S8_Dataset}
{\bf Hostplants of Nymphalidae
butterflies recorded from the literature}.

\subsection*{S9 Dataset}
\label{S9_Dataset}
{\bf References for hostplants data}.

\subsection*{S10 Supporting Information}
\label{S10_Supporting_Information}
{\bf R script to run MEDUSA on the MCC Nymphalidae tree from
\cite{wahlberg2009}}. This script removes the outgroup
taxa and loads the richness data for the tree terminals.

\subsection*{S11 Supporting Information}
\label{S11_Supporting_Information}
{\bf R script to run a MultiMEDUSA analysis on 1000 random trees from
\cite{wahlberg2009}}.

\subsection*{S12 Supporting Information}
\label{S12_Supporting_Information}
{\bf R script to run a MEDUSA
analysis on the MCC tree from the 1000 random trees selected from the
posterior distribution}.

\subsection*{S13 Dataset}
\label{S13_Dataset}
{\bf Data matrix with character states for hosplant use}.

\subsection*{S14 Supporting Information}
\label{S14_Supporting_Information}
{\bf MCC tree from the 1000 random
trees selected from the posterior distribution}.

\subsection*{S15 File}
\label{S15_File}
{\bf Raw results from the MultiMEDUSA run on the random sample of trees from 
the posterior distribution}.

\subsection*{S16 Text}
\label{S16_Text}
{\bf Summary results from a
MultiMEDUSA run on the 1000 random trees selected from the posterior
distribution}.

\subsection*{S17 Figure}
\label{S17_Figure}
{\bf Probability of nodes with
estimated rates from a MultiMEDUSA run on the 1000 random trees selected
from the posterior distribution}.

\subsection*{S18 Supporting Information}
\label{S18_Supporting_Information}
{\bf R script for running the BiSSE analysis}.

\subsection*{S19 Supporting Information}
\label{S19_Supporting_Information}
{\bf Raw results for the BiSSE analysis}.

\subsection*{S20 Supporting Information}
\label{S20_Supporting_Information}
{\bf Raw results for the combined BiSSE analysis}.

\subsection*{S21 Dataset}
\label{S21_Dataset}
{\bf MCC Nymphalidae tree from \cite{wahlberg2009}}

\subsection*{S22 Text}
\label{S22_Text}
{\bf Results of running MEDUSA on the MCC tree}.

\subsection*{S23 Figure}
\label{S23_Figure}
{\bf Summary of the MultiMEDUSA
analysis on the 1000 trees from the posterior distribution}.

\subsection*{S24 Text}
\label{S24_Text}
{\bf Summary results from a MEDUSA
run on the MCC tree from the 1000 random trees selected from the
posterior distribution}.

\subsection*{S25 Figure}
\label{S25_Figure}
{\bf Figure for MEDUSA run on MCC tree from random 1000 trees}.

\section*{Acknowledgments}
We are thankful to Mark Cornwall for help with the script to extend
MEDUSA to include phylogenetic uncertainty, Niklas Wahlberg for
commenting on the manuscript and giving us the posterior distribution of
trees, Luke Harmon for commenting on the manuscript, Jeffrey C. Oliver and anonymous
reviewers for their comments, which greatly improved the manuscript,
Jessica Slove Davidson and Niklas Janz for access to their hostplant
data. The study was supported by a Kone Foundation grant (awarded to
Niklas Wahlberg), Finland (C. Pe\~na) and the Research Council of Norway
(grant no. 204308 to M. Espeland). We acknowledge CSC--IT Center for
Science Ltd. (Finland) for the allocation of computational resources.

\nolinenumbers


\bibliography{refs}{}


\end{document}
